\begin{proposition}
The equation
\[
    ax_1^2+bx_2^2+cx_3^2=0,\quad abc\ne0,\quad (x_1,x_2,x_3)=1
\]
has only the trivial solution $x_1=0,\ x_2=0,\ x_3=0$ if either
\[
    a\equiv b\equiv c\equiv1\pmod2,\qquad a\equiv b\equiv c\pmod4,
\]
or
\[
    a/2\equiv b\equiv c\equiv1\pmod2,\qquad b+c\equiv a\text{ or }4\pmod8.
\]
\end{proposition}

\begin{proposition}
Assume that $a,b,c>0$ are square-free and relatively prime in pairs, and that
\[
    ax^2+by^2-cz^2\equiv0\pmod{p^r}
\]
is solvable for all primes $p$ and integers $r\ge1$ with $(x,y,z,p)=1$. Then
\[
    ax^2+by^2-cz^2=0
\]
has a non-trivial solution with
\[
    |x|<\sqrt{2bc},\qquad |y|<\sqrt{2ca},\qquad |z|<2\sqrt{ab}.
\]
\end{proposition}

\begin{proposition}
If $a,b,c$ are square-free, $(a,b)=(b,c)=(c,a)=1$, and $a,b,c$ do not have
the same sign, then the equation
\[
    ax^2+by^2+cz^2=0
\]
has non-trivial integer solutions if and only if $-bc$ is a quadratic residue
of $a$, and so for every prime factor of $a$, i.e. $x^2+bc\equiv0\pmod a$ is
solvable, and similarly for $-ca$ and $-ab$, and if
\[
    ax^2+by^2+cz^2\equiv0\pmod8
\]
is solvable. These give $N$, say, necessary and sufficient conditions. If any
$N-1$ of these conditions are satisfied, so is the remaining one.
\end{proposition}

\begin{proposition}
The solvable equation
\[
    ax^2+by^2+cz^2=0
\]
taken in the canonical form with $a>0$, $b>0$, $c<0$ has a non-trivial solution
with
\[
    |x|\le\sqrt{b|c|},\qquad |y|\le\sqrt{|c|a},\qquad |z|\le\sqrt{ab}.
\]
All the equality signs may be removed unless two of the $a,b,c$ are unity.
\end{proposition}
